\section{Introduction}

\neuroSymbolicAI{} can be approached in the tensor network formalism \tnreason{} \cite{goessmann_tensor_2026}, unifying logical, neural and probabilistic paradigms of artificial intelligence.
Quantum Circuits are also, by the axioms of quantum mechanics, tensor networks.
This motivates this work, which investigates quantum algorithms for inference leading to efficiency increases compared with classical alternatives.
We study an rejection sampling inference scheme of the broad model class.
Quantum rejection sampling then brings an square root quantum advantage.
To be more precise, we consider \ComputationActivationNetworks{}, a generic model class in \tnreason{} fo \neuroSymbolicAI{}.

\subsection{Related works}

Quantum rejection sampling has been developed in these works:
\begin{itemize}
    \item \cite{grover_fast_1996} Grover algorithm (search in unstructured database)
    \item \cite{brassard_exact_1997,brassard_quantum_2002} generalization to amplitude amplification
    \item \cite{ozols_quantum_2013} introduced quantum rejection sampling (using amplitude amplification)
    \item \cite{low_quantum_2014} used quantum rejection sampling for Bayesian network sampling (which is NP-hard when conditioned on evidence, see e.g. \cite{koller_probabilistic_2009})
    \item No exponential speedups expected by Quantum Computing in AI \cite{bennett_strengths_1997}
\end{itemize}

Several applications of Quantum Inference on symbolic AI have been studied.
The novelty of our approach is a concise framework extending these approaches and connecting with generic neural decompositions of functions (see \computationCircuits{}).

Main approach for sampling: Quantum Inference scheme on Bayesian networks \cite{low_quantum_2014}
\begin{itemize}
    \item Extend to more general tensor networks than Bayesian networks: \ComputationActivationNetworks{}
\end{itemize}

Further literature:
\begin{itemize}
    \item \cite{schuld_machine_2021}: Sect 7.3.2 Reviewing the paper \cite{low_quantum_2014} as an application of fault-tolerant quantum computing
    \item \cite{wittek_quantum_2017}: Review of Markov Logic Network sampling (which are a special case of \ComputationActivationNetworks{})
\end{itemize}


\subsection{Contributions}

We apply an adjusted quantum rejection sampling approach of \cite{low_quantum_2014} for \ComputationActivationNetworks{} \cite{goessmann_tensor_2026}.
The main adjustments are:
\begin{itemize}
    \item {\bf Sparse representation of uniformly controlled unitaries:} We apply similar sparsity mechanisms as for classical \ComputationActivationNetworks{} to decompose the uniformly controlled unitaries.
    \item {\bf Ancilla augmentation:} To sample also from generic tensor networks, we introduce ancilla variables such that generic tensor networks are conditioned Bayesian networks.
\end{itemize}